\section{Exact algebra of the new period-six orbit}

The new word \(B_6=000231\) has sign coordinates \(---++-\).  Set
\begin{equation}\label{eq:b6-coordinates}
 c=-\frac{\sqrt7}{6},\qquad
 D=\sqrt{25+4\sqrt7},\qquad
 a=\frac{-1-D}{12},\qquad d=\frac{-1+D}{12}.
\end{equation}

\begin{proposition}[radical period-six certificate]\label{prop:b6}
The cyclic coordinate tuple
\[
 (q_0,\ldots,q_5)=(a,a,c,d,d,c)
\]
satisfies
\(q_{j+1}=1-6q_j^2-q_{j-1}\) exactly and has sign word \(---++-\).
Its monodromy trace and trace polynomial are
\begin{align}
 T_6&=18062+5352\sqrt7,\label{eq:b6-trace}\\
 Q_6(T)&=T^2-36124T+125728516.\label{eq:b6-trace-poly}
\end{align}
The signed unstable multiplier has irreducible reciprocal polynomial
\begin{equation}\label{eq:b6-multiplier}
 z^4-36124z^3+125728518z^2-36124z+1.
\end{equation}
Its Galois excess is
\begin{equation}\label{eq:b6-excess}
 \Egal_{B_6}=\acosh(9031-2676\sqrt7).
\end{equation}
\end{proposition}

\begin{proof}
Substitution of \eqref{eq:b6-coordinates} into the six cyclic recurrence
equations gives zero.  Multiplying the derivative matrices
\[
 DH_6(q_j,q_{j-1})=\begin{pmatrix}-12q_j&-1\\1&0\end{pmatrix}
\]
gives determinant one and trace \eqref{eq:b6-trace}.  Conjugating
\(\sqrt7\mapsto-\sqrt7\) yields \eqref{eq:b6-trace-poly}.  Eliminating
\(T=z+z^{-1}\) gives \eqref{eq:b6-multiplier}; direct factorization over
\(\mathbb Q\) is irreducible.  Concretely, modulo \(13\) it becomes
\[
 \bar f(z)=z^4+3z^3+6z^2+3z+1,
\]
and
\(\gcd(\bar f,z^{13}-z)=\gcd(\bar f,z^{169}-z)=1\).
Thus it has no linear or quadratic factor over \(\mathbb F_{13}\), proving
irreducibility over \(\mathbb Q\).  The nonphysical trace is
\(18062-5352\sqrt7\), so its expanding reciprocal pair contributes
\(\acosh(|T'|/2)\), which is \eqref{eq:b6-excess}.
\end{proof}

The conjugate half-trace is isolated without floating point:
\begin{equation}\label{eq:b6-isolation}
 1950<9031-2676\sqrt7<1951.
\end{equation}
Indeed
\[
 7081^2-7\cdot2676^2=13729>0,
 \qquad
 7\cdot2676^2-7080^2=432>0.
\]
Numerically, \(\Egal_{B_6}=8.26922881806116\ldots\); decimals are not used
in the obstruction proof.
